\chapter{Dynamik: Kreislauf und Grenzzyklus}

Die fraktale Energieverteilung ist nicht statisch – sie ist der \emph{zeitliche Mittelwert} eines permanenten, geschlossenen Kreislaufs.

\subsection{Der geschlossene Energiekreislauf}

Die Energie zirkuliert in drei Phasen:

\begin{equation}
\text{Vakuum} \xrightarrow{\sigma_{\text{cre}}} \text{Materie} \xrightarrow{\sigma_{\text{drift}}} \Omega\text{-Feld} \xrightarrow{\sigma_{\text{back}}} \text{Vakuum}
\end{equation}

Die Flussraten sind:

\begin{align}
\sigma_{\text{cre}} &= \frac{1}{T} \int \frac{|\nabla Q|^2}{\hbar c} \, d^3r > 0, \\
\sigma_{\text{drift}} &= \frac{1}{T} \int \frac{\dot{\beta}_{\text{DSTT}}}{\beta_{\text{DSTT}}} \, d^3r > 0, \\
\sigma_{\text{back}} &= \frac{1}{T} \int \frac{|\vec{\Omega}|^2}{8\pi G} \, d^3r > 0.
\end{align}

Im stationären Zustand gilt:

\begin{equation}
\sigma_{\text{cre}} = \sigma_{\text{drift}} = \sigma_{\text{back}} \equiv \sigma.
\end{equation}

\subsection{Der Grenzzyklus}

Die globale Entropie ist konstant – das ist das Resultat eines \emph{attraktiven Grenzzyklus}. Es gibt keinen Wärmetod, weil:

\begin{itemize}
    \item Die fraktale Entropie \(S_D\) kein globales Maximum hat (\(D < 3\)).
    \item Die DSTT-Drift (\( \dot{\beta} > 0 \)) permanente radiale Energiegradienten erzeugt.
    \item Die Flussraten \( \sigma_{\text{cre}}, \sigma_{\text{drift}}, \sigma_{\text{back}} \) immer positiv sind.
\end{itemize}

Die DSTT-Drift ist der Motor des Kreislaufs – sie ist die Quelle der Irreversibilität und des intrinsischen Zeitpfeils.